\chapter{Kết luận và hướng phát triển}

\section{Kết luận}

Đồ án đã xây dựng thành công hệ thống nhúng hỗ trợ bán hàng thông minh
ShopNexus, bao gồm một pipeline xử lý AI hoàn chỉnh từ thu thập dữ liệu
đến hiển thị kết quả. Các kết quả chính đạt được:

\textbf{Về thiết bị nhúng,} ESP32-CAM khai thác hiệu quả kiến trúc dual-core
để thu đồng thời âm thanh 16 kHz và video JPEG mà không xảy ra hiện tượng
buffer underrun. Giao thức điều khiển camera tự thiết kế cho phép thay đổi
21 tham số real-time từ giao diện web.

\textbf{Về khử nhiễu,} mô hình DTLN được huấn luyện trên dữ liệu tiếng Việt
VIVOS cho kết quả khử nhiễu tốt với chỉ 1,8 triệu tham số (3,9 MB) và
latency 32 ms, phù hợp yêu cầu real-time. Kiến trúc hai giai đoạn
(STFT + Conv1D Encoder) khử nhiễu toàn diện hơn phương pháp đơn miền.

\textbf{Về nhận dạng và sinh mô tả,} dịch vụ Voice Rephraze kết hợp
gpt-4o-transcribe (STT) và GPT-4o-mini (sinh mô tả) tạo ra mô tả sản phẩm
chất lượng với 4+ phong cách viết. Prompt engineering xử lý hiệu quả text
lộn xộn từ STT, bỏ qua từ lấp và lỗi phát âm.

\textbf{Về phân loại sản phẩm,} XLM-RoBERTa đạt accuracy 90,14\% và F1-score
90,00\% trên 32 danh mục thương mại điện tử, với thời gian suy luận 45 ms
trên CPU --- đáp ứng yêu cầu ứng dụng thực tế.

Nhìn tổng thể, hệ thống đã chứng minh tính khả thi của việc kết hợp thiết bị
nhúng giá rẻ (ESP32, giá khoảng 5 USD) với các dịch vụ AI hiện đại để tạo ra
giải pháp hỗ trợ bán hàng tự động. Người bán chỉ cần mô tả sản phẩm bằng
giọng nói, hệ thống tự động xử lý toàn bộ quy trình còn lại --- tiết kiệm
thời gian và nâng cao chất lượng mô tả sản phẩm.

\section{Hướng phát triển}

Dựa trên kết quả và hạn chế đã phân tích, nhóm đề xuất các hướng phát triển
sau cho hệ thống ShopNexus:

\begin{itemize}
\item \textbf{Triển khai PhoWhisper local:} Thay thế cloud API OpenAI bằng mô
  hình PhoWhisper (VinAI) chạy trực tiếp trên server, giảm phụ thuộc internet
  và chi phí vận hành. PhoWhisper được huấn luyện đặc biệt cho tiếng Việt nên
  có thể cho kết quả tương đương hoặc tốt hơn trong ngữ cảnh Việt Nam.

\item \textbf{Bổ sung multi-microphone beamforming:} Sử dụng 2--4 microphone
  INMP441 kết hợp thuật toán beamforming (delay-and-sum hoặc MVDR) để tăng
  SNR đầu vào trước khi đưa vào DTLN, cải thiện đáng kể chất lượng trong
  môi trường nhiều nhiễu.

\item \textbf{Cải tiến DTLN:} Thêm windowing function (Hann hoặc Hamming)
  trong bước STFT để giảm spectral leakage. Nghiên cứu ước lượng pha sạch
  thay vì tái sử dụng pha gốc ở giai đoạn 1, giảm artifact âm thanh.

\item \textbf{Tối ưu phân loại cho tiếng Việt:} Fine-tune thêm mô hình
  XLM-RoBERTa trên dữ liệu mô tả sản phẩm tiếng Việt thực tế từ các
  sàn thương mại điện tử Việt Nam (Shopee, Lazada), nâng cao accuracy
  cho các danh mục đặc thù thị trường nội địa.

\item \textbf{Triển khai edge inference:} Chuyển DTLN sang TensorFlow Lite
  và triển khai trực tiếp trên Raspberry Pi 4 hoặc thiết bị edge tương
  đương, giảm latency mạng và cho phép hoạt động semi-offline.

\item \textbf{Buffer offline cho ESP32:} Bổ sung cơ chế lưu trữ tạm audio
  vào thẻ microSD khi mất kết nối WiFi, tự động gửi lại khi kết nối được
  khôi phục --- đảm bảo không mất dữ liệu trong môi trường mạng không ổn định.

\item \textbf{Hỗ trợ đa nền tảng:} Phát triển ứng dụng mobile (Android/iOS)
  song song với giao diện web, cho phép người bán sử dụng microphone trên
  điện thoại khi không có thiết bị ESP32.
\end{itemize}
